
\documentclass[11pt,a4paper]{article}

\usepackage[english]{babel}
\usepackage[utf8]{inputenc}
\usepackage[T1]{fontenc}
\usepackage{lmodern}
\usepackage{geometry}
\usepackage{graphicx}
\usepackage{booktabs}
\usepackage{tabularx}
\usepackage{array}
\usepackage{enumitem}
\usepackage{parskip}
\usepackage{microtype}
\usepackage[hidelinks]{hyperref}
\usepackage[most]{tcolorbox}
\usepackage{xcolor}
\usepackage{fancyhdr}
\usepackage{titlesec}
\usepackage{float}

\geometry{margin=22mm}

\definecolor{ManualBlue}{HTML}{2F5D7C}
\definecolor{ManualGray}{HTML}{F4F5F6}
\definecolor{ManualDark}{HTML}{222222}
\definecolor{WarningRed}{HTML}{A43B32}
\definecolor{NoteGreen}{HTML}{3B6E52}

\hypersetup{
    colorlinks=true,
    linkcolor=ManualBlue,
    urlcolor=ManualBlue
}

\titleformat{\section}
  {\Large\bfseries\color{ManualDark}}
  {\thesection}{0.75em}{}

\titleformat{\subsection}
  {\large\bfseries\color{ManualDark}}
  {\thesubsection}{0.75em}{}

\setlist[itemize]{leftmargin=1.6em,itemsep=0.25em,topsep=0.35em}
\setlist[enumerate]{leftmargin=1.8em,itemsep=0.35em,topsep=0.35em}

\newtcolorbox{importantbox}{
  colback=ManualGray,
  colframe=ManualBlue,
  boxrule=0.7pt,
  arc=2mm,
  left=2.5mm,right=2.5mm,top=2mm,bottom=2mm
}

\newtcolorbox{warningbox}{
  colback=red!3,
  colframe=WarningRed,
  boxrule=0.8pt,
  arc=2mm,
  left=2.5mm,right=2.5mm,top=2mm,bottom=2mm
}

\newtcolorbox{tipbox}{
  colback=green!3,
  colframe=NoteGreen,
  boxrule=0.7pt,
  arc=2mm,
  left=2.5mm,right=2.5mm,top=2mm,bottom=2mm
}

\pagestyle{fancy}
\fancyhf{}
\lhead{Rockit Rev. 1}
\rhead{User Manual -- Firmware 3.0}
\cfoot{\thepage}

\renewcommand{\arraystretch}{1.2}

\begin{document}

\begin{titlepage}
\begin{center}

\vspace*{12mm}

{\Huge\bfseries Rockit Rev. 1}\\[3mm]
{\LARGE User Manual}\\[4mm]
{\Large Firmware 3.0}\\[15mm]

\includegraphics[width=0.78\textwidth]{JSTconnector.jpg}

\vfill

{\Large\bfseries Triskel Instruments}\\[2mm]
{\normalsize formerly DanInvents}\\[8mm]

{\large \today}

\end{center}
\end{titlepage}

\tableofcontents
\newpage

\section{About this manual}

This manual explains how to prepare, configure, fly and update the Rockit Rev. 1 flight computer using Firmware 3.0.

Rockit was originally released under the name \textbf{DanInvents}. The product is now developed and supported by \textbf{Triskel Instruments}. Older boards, manuals or videos may still carry the DanInvents name.

\begin{importantbox}
\textbf{Firmware 3.0 is a major software update.} It introduces a new MicroPython-based firmware, improved launch and apogee detection, safer recovery sequencing, persistent settings, and a standalone Windows updater.
\end{importantbox}

\section{Disclaimer}

In no respect shall Triskel Instruments be accountable for liabilities, claims, demands, damages or suits resulting from the use of Rockit and/or its related firmware.

By using this product, you accept the risks associated with model rocketry, lithium batteries, servo-driven recovery systems and experimental electronics.

Always test your complete recovery system on the ground before flight.

\section{What is Rockit?}

Rockit, or \textit{Rocket Operation Computing Kit}, is a compact flight computer designed for model rocketry.

It provides:

\begin{itemize}
    \item altitude and acceleration measurement;
    \item automatic or timed recovery deployment;
    \item control of two servo motors;
    \item flight-data logging to a microSD card;
    \item maximum-altitude reporting after flight;
    \item battery charging for a single-cell lithium battery.
\end{itemize}

\subsection{Main specifications}

\begin{table}[H]
\centering
\begin{tabularx}{\textwidth}{>{\bfseries}l X}
\toprule
Specification & Value \\
\midrule
Dimensions & 44 mm $\times$ 22 mm $\times$ 9 mm \\
Mass & 5 g \\
Input voltage & 2.3 to 5.6 V \\
Idle current & Approx. 70 mA \\
Logging current & Approx. 80 mA \\
Peak current while beeping & Approx. 200 mA \\
Typical current while moving a servo & Approx. 400 mA \\
Maximum continuous 5 V output current & 2 A \\
Battery charging current & 500 mA \\
Sampling rate & 100 Hz \\
Altimeter range & Up to 9160 m \\
Accelerometer range & Up to 22.6 g \\
\bottomrule
\end{tabularx}
\end{table}

\begin{tipbox}
The practical lower voltage limit is often set by the protection circuit inside the LiPo battery. Many protected batteries disconnect themselves at around 3 V.
\end{tipbox}

\section{Board layout}

\begin{figure}[H]
\centering
\includegraphics[width=0.95\textwidth]{Rockit dimensions.png}
\caption{Rockit Rev. 1 dimensions. The 2.54 mm holes are connected to ground and can be used for probing.}
\end{figure}

\begin{figure}[H]
\centering
\includegraphics[width=0.65\textwidth]{rockitLayout.png}
\caption{Top-side layout. The I\textsuperscript{2}C pins and servo signal pins use 3.3 V logic.}
\end{figure}

\begin{warningbox}
\textbf{Do not apply 5 V logic to the 3.3 V signal pins.} Doing so may damage the RP2040 microcontroller.
\end{warningbox}

\begin{figure}[H]
\centering
\includegraphics[width=0.65\textwidth]{bottomLayout.png}
\caption{Bottom-side layout. Servo 1 uses GPIO 28 and Servo 2 uses GPIO 27. I\textsuperscript{2}C SDA and SCL use GPIO 10 and GPIO 11.}
\end{figure}

\section{Before the first flight}

Before flying Rockit, check the following:

\begin{enumerate}
    \item Insert a compatible microSD card.
    \item Connect the battery with the correct polarity.
    \item Configure both servo start and end positions.
    \item Configure the required deployment delays.
    \item Run the position-3 recovery-system test.
    \item Confirm that the battery can move both servos without shutting down.
    \item Make sure the pressure sensor is protected from direct sunlight but its vent hole is not blocked.
\end{enumerate}

\section{Power and battery}

Rockit is designed to operate from 2.3 to 5.6 V.

The onboard battery charger is intended only for a \textbf{single-cell (1S) lithium-ion or lithium-polymer battery}, typically 3.7 V nominal.

\begin{warningbox}
\textbf{Battery polarity is critical.} If Rockit becomes unusually hot immediately after connecting the battery, disconnect it at once.
\end{warningbox}

\begin{figure}[H]
\centering
\includegraphics[width=0.48\textwidth]{JSTconnector.png}
\caption{Correct JST battery polarity. Red is battery positive and black is ground.}
\end{figure}

Servo motors can draw large current spikes. A battery with an aggressive protection circuit may shut itself off when a servo moves.

\textbf{Always test the complete battery and servo combination before flight.}

\section{Choosing a microSD card}

For best compatibility:

\begin{itemize}
    \item use a Class 10 microSD card;
    \item use 16 GB or less when possible;
    \item format larger cards as FAT32 if necessary.
\end{itemize}

If Rockit appears to start but does not proceed normally, the microSD card is one of the first things to check.

\section{Rotary switch quick reference}

The rotary switch is used both to select the flight mode and to configure Rockit.

\begin{table}[H]
\centering
\begin{tabularx}{\textwidth}{>{\centering\arraybackslash\bfseries}p{1.4cm} X}
\toprule
Position & Function \\
\midrule
A & Automatic mode: deploy at detected apogee \\
B & Timer mode: deploy after a preset time from launch detection \\
C & Set Servo 1 delay after apogee in Automatic mode \\
D & Set Servo 1 delay after launch in Timer mode \\
E & Set Servo 1 start position \\
F & Set Servo 1 final position \\
0 & Set Servo 2 start position \\
1 & Set Servo 2 final position \\
2 & Set additional delay before Servo 2 moves \\
3 & Test the complete servo/recovery sequence \\
4--9 & Inactive / non-flight positions \\
\bottomrule
\end{tabularx}
\end{table}

\begin{importantbox}
Some Rockit boards use a standard-coded rotary switch and some use a complementary-coded switch. Firmware 3.0 supports both. The Windows updater can detect the correct type automatically.
\end{importantbox}

\section{Changing settings}

The basic configuration workflow is:

\begin{enumerate}
    \item Turn Rockit off.
    \item Turn the rotary switch to the setting you want to change.
    \item Turn Rockit on.
    \item Follow the LED/buzzer feedback.
    \item Rotate the switch to choose the desired value.
    \item Keep the value stable so Rockit can save it.
    \item Turn Rockit off when the setting has been stored.
\end{enumerate}

Firmware 3.0 stores the settings in a persistent configuration file, so they remain available after power is removed and can also be backed up by the firmware updater.

\subsection{Setting the deployment delays}

There are three timing settings:

\begin{itemize}
    \item \textbf{Position C}: delay between apogee detection and Servo 1 movement in Automatic mode;
    \item \textbf{Position D}: delay between launch detection and Servo 1 movement in Timer mode;
    \item \textbf{Position 2}: additional delay before Servo 2 moves.
\end{itemize}

\subsubsection*{Automatic mode delay -- position C}

Each rotary step adds \textbf{0.5 seconds}.

\begin{center}
\begin{tabular}{c|cccccccc}
Switch value & 0 & 1 & 2 & 3 & 4 & 5 & 10 & 15 \\
\hline
Delay (s) & 0 & 0.5 & 1.0 & 1.5 & 2.0 & 2.5 & 5.0 & 7.5
\end{tabular}
\end{center}

\subsubsection*{Timer mode delay -- position D}

Each rotary step adds \textbf{1 second}.

\begin{center}
\begin{tabular}{c|cccccccc}
Switch value & 0 & 1 & 2 & 3 & 4 & 5 & 10 & 15 \\
\hline
Delay (s) & 0 & 1 & 2 & 3 & 4 & 5 & 10 & 15
\end{tabular}
\end{center}

\subsubsection*{Servo 2 additional delay -- position 2}

Each rotary step adds \textbf{0.5 seconds} after Servo 1.

\begin{center}
\begin{tabular}{c|cccccccc}
Switch value & 0 & 1 & 2 & 3 & 4 & 5 & 10 & 15 \\
\hline
Additional delay (s) & 0 & 0.5 & 1.0 & 1.5 & 2.0 & 2.5 & 5.0 & 7.5
\end{tabular}
\end{center}

\begin{tipbox}
For Servo 2, think of the setting as ``how long after Servo 1 should Servo 2 move?'' This is usually easier to understand than calculating the total time from launch or apogee.
\end{tipbox}

\subsection{Setting the servo positions}

Use:

\begin{itemize}
    \item \textbf{E} for Servo 1 start position;
    \item \textbf{F} for Servo 1 final position;
    \item \textbf{0} for Servo 2 start position;
    \item \textbf{1} for Servo 2 final position.
\end{itemize}

When adjusting a servo, the rotary switch gives 16 possible positions from 0 to F.

You do not need to calculate an angle. Simply choose the position that places your mechanism where you want it.

\begin{tipbox}
If a servo does not appear to move during deployment, check that its start and final positions are not set to the same value.
\end{tipbox}

\section{Testing the recovery system}

Position \textbf{3} performs a complete ground test of the recovery servos.

With Rockit turned off:

\begin{enumerate}
    \item turn the rotary switch to \textbf{3};
    \item power Rockit on;
    \item Servo 1 moves to its programmed start position;
    \item Servo 2 moves to its programmed start position;
    \item Rockit waits approximately 5 seconds;
    \item Servo 1 moves to its programmed final position;
    \item Servo 2 moves to its programmed final position;
    \item the flight computer then remains inactive.
\end{enumerate}

\begin{importantbox}
Use position 3 before every flight in which you have changed the recovery mechanism, servo linkage, battery, or servo settings.
\end{importantbox}

\section{Flight modes}

\subsection{Automatic mode -- position A}

Automatic mode is intended for normal recovery deployment at apogee.

After startup:

\begin{enumerate}
    \item Rockit arms after approximately 3 seconds.
    \item Rockit waits for launch.
    \item Launch is confirmed using both altitude and acceleration information.
    \item During ascent, Rockit tracks the filtered altitude.
    \item After apogee is confirmed, Rockit waits for the delay configured with position C.
    \item Servo 1 moves.
    \item Servo 2 moves after the additional delay configured with position 2.
\end{enumerate}

\subsection{How launch detection works}

Rockit does not trigger from a single bump or a single pressure change.

For a real flight, the firmware looks for:

\begin{itemize}
    \item an altitude increase of about \textbf{10 m};
    \item a change in acceleration of more than about \textbf{2 g};
    \item the acceleration condition to persist for multiple samples;
    \item both events to occur close together in time.
\end{itemize}

This makes false launch detection from handling or vibration much less likely than in earlier firmware versions.

\subsection{How apogee detection works}

Firmware 3.0 no longer declares apogee just because one or two altitude samples decrease.

Instead, it:

\begin{itemize}
    \item tracks the highest filtered altitude reached;
    \item ignores very early false peaks immediately after launch;
    \item waits until the rocket has climbed to a meaningful altitude;
    \item requires a clear drop from the peak;
    \item requires that drop to persist briefly before deployment begins.
\end{itemize}

\begin{importantbox}
Rockit waits for a clear and persistent descent from the detected altitude peak before treating the rocket as being past apogee.
\end{importantbox}

\subsection{Timer mode -- position B}

Timer mode deploys the recovery system after a fixed time from launch detection.

After launch is confirmed:

\begin{enumerate}
    \item Rockit waits for the time configured with position D.
    \item Servo 1 moves.
    \item Rockit waits for the additional Servo 2 delay configured with position 2.
    \item Servo 2 moves.
\end{enumerate}

Timer mode is useful when you want predictable deployment timing rather than apogee-based deployment.

\section{What happens after launch?}

During flight, Rockit records data to the microSD card.

The logged data include:

\begin{itemize}
    \item time;
    \item altitude;
    \item filtered altitude;
    \item longitudinal acceleration;
    \item perpendicular acceleration;
    \item temperature.
\end{itemize}

Rockit also keeps a short pre-launch buffer, so the saved file contains data from immediately before takeoff.

Approximately one minute after launch, Rockit transitions to maximum-altitude reporting.

\section{Reading the maximum altitude}

After the flight, Rockit repeatedly reports the maximum altitude using the buzzer.

The pattern is:

\begin{enumerate}
    \item two short beeps mark the start;
    \item each altitude digit is beeped separately;
    \item two short beeps mark the end;
    \item the pattern repeats.
\end{enumerate}

The digit \textbf{0} is represented by a long beep.

\subsection*{Example: 67 m}

You will hear:

\begin{center}
two short beeps $\rightarrow$ six short beeps $\rightarrow$ seven short beeps $\rightarrow$ two short beeps
\end{center}

The complete flight record can also be read directly from the CSV file stored on the microSD card.

\section{Updating Rockit}

Firmware 3.0 introduces a standalone Windows updater.

You no longer need Python, Arduino IDE, VS Code or \texttt{mpremote}.

\subsection{Recommended update method}

\begin{enumerate}
    \item Download the latest Rockit firmware updater from:
    \begin{center}
    \url{https://github.com/DanInvents/Rockit/releases}
    \end{center}
    \item Connect Rockit to the computer using USB.
    \item Run \texttt{Rockit\_Firmware\_Updater.exe}.
    \item Follow the instructions shown on screen.
    \item If asked, turn the rotary switch to position A so the updater can identify the switch coding.
    \item Wait until the updater displays \textbf{UPDATE SUCCESSFUL}.
\end{enumerate}

The updater can:

\begin{itemize}
    \item identify an existing MicroPython installation;
    \item back up saved Rockit settings;
    \item install Firmware 3.0;
    \item restore servo positions and recovery delays;
    \item detect standard or complementary-coded rotary switches;
    \item guide older Arduino-based Rockit units through BOOTSEL migration;
    \item install the validated MicroPython runtime automatically.
\end{itemize}

\begin{importantbox}
If you are updating an older Arduino-based Rockit, its old EEPROM settings cannot be transferred automatically. Configure and test the servos again after the migration.
\end{importantbox}

\section{What changed in Firmware 3.0?}

Firmware 3.0 is the largest software revision since Rockit was introduced.

The most important improvements are:

\begin{itemize}
    \item more reliable launch detection using altitude and acceleration together;
    \item more robust automatic apogee detection;
    \item improved recovery-servo timing and sequencing;
    \item improved servo startup and detach behavior;
    \item a dedicated position-3 recovery-system test;
    \item a shorter arming delay;
    \item persistent human-readable settings;
    \item settings backup and restore during firmware updates;
    \item support for both rotary-switch hardware variants;
    \item improved pre-launch buffering and SD-card logging;
    \item clearer maximum-altitude reporting;
    \item a much faster post-flight report;
    \item improved buzzer and PWM handling;
    \item standalone Windows firmware updates;
    \item migration from the previous Arduino firmware to MicroPython.
\end{itemize}

Firmware 3.0 has been flight-tested successfully in both:

\begin{itemize}
    \item \textbf{A -- Automatic mode};
    \item \textbf{B -- Timer mode}.
\end{itemize}

\section{Frequently asked questions}

\subsection*{My flight computer gets very hot when I connect the battery. What should I do?}

Disconnect the battery immediately. The battery connector polarity is probably reversed.

Do not leave the battery connected while diagnosing the problem.

\subsection*{My servo does not move.}

Check the following:

\begin{itemize}
    \item the servo is connected to the correct output;
    \item the battery can supply enough current;
    \item the start and final positions are different;
    \item the mechanism is not mechanically jammed;
    \item the position-3 ground test works correctly.
\end{itemize}

\subsection*{Rockit does not start normally.}

Check the microSD card first.

For best compatibility, use a Class 10 card of 16 GB or less.

Also check that the battery voltage is adequate.

\subsection*{Can I test launch detection by tilting the board?}

No. Firmware 3.0 uses both acceleration and altitude changes to detect launch.

For checking the recovery mechanism, use rotary position \textbf{3} instead of trying to simulate a complete launch.

\subsection*{What is the difference between standard and complementary-coded Rockit boards?}

Some production batches used rotary switches with opposite electrical coding.

Firmware 3.0 supports both types using the same firmware.

The Windows updater can detect the correct type automatically by asking you to place the switch at position \textbf{A}.

\subsection*{Do I need Arduino IDE or Python to update Rockit?}

No.

For Firmware 3.0, use the standalone Windows updater from the Rockit GitHub Releases page.

\section{Support and firmware downloads}

Firmware releases and update information are available at:

\begin{center}
\url{https://github.com/DanInvents/Rockit/releases}
\end{center}

Rockit is developed and supported by:

\begin{center}
\textbf{Triskel Instruments}\\
formerly DanInvents
\end{center}

\end{document}